\FigureLayoutDeclare{img/2014/hunan_liberal/q19_fig1.png}{12.0cm}{10bp 0bp 42bp 24bp}
\FigureLayoutDeclare{img/2014/hunan_liberal/q20_fig1.png}{7.0cm}{14bp 11bp 22bp 4bp}

\chapter{2014年湖南卷（文）}

\section{单选题}

\begin{problem}
设命题 \(p: \forall x \in \R\)，\(x^2 + 1 > 0\)，则 \(\neg p\) 为（\quad）

\choices
    {\(\exists x_0 \in \R, x_0^2 + 1 > 0\)}
    {\(\exists x_0 \in \R\)，\(x_0^2 + 1 \leqslant 0\)}
    {\(\exists x_0 \in \R\)，\(x_0^2 + 1 < 0\)}
    {\(\forall x \in \R\)，\(x^2 + 1 \leqslant 0\)}
\end{problem}

\begin{answer}
B
\end{answer}

\begin{solution}
由全称命题的否定规则，\(\neg(\forall x\in\R, P(x))\) 等价于 \(\exists x_0\in\R\)，使 \(\neg P(x_0)\) 成立. 因此 \(\neg p\) 为 \(\exists x_0\in\R\)，\(x_0^2+1\leqslant0\)，故选 B.
\end{solution}

\begin{problem}
已知集合 \(A = \{x \mid x > 2\}\)，\(B = \{x \mid 1 < x < 3\}\)，则 \(A \cap B =\)（\quad）

\choices
    {\( \{x \mid x > 2\} \)}
    {\( \{x \mid x > 1\} \)}
    {\(\{x\mid 2 < x < 3\}\)}
    {\(\{x\mid 1 < x < 3\}\)}
\end{problem}

\begin{answer}
C
\end{answer}

\begin{solution}
由交集的定义，\(x\in A\cap B\) 当且仅当 \(x>2\) 且 \(1<x<3\)，所以 \(A\cap B=\{x\mid 2<x<3\}\)，故选 C.
\end{solution}

\begin{problem}
对一个容量为 \(N\) 的总体抽取容量为 \(n\) 的样本，当选取简单随机抽样，系统抽样和分层抽样三种不同方法抽取样本时，总体中每个个体被抽中的概率分别为 \(p_1\)，\(p_2\)，\(p_3\)，则（\quad）

\choices
    {\(p_1 = p_2 < p_3\)}
    {\(p_2 = p_3 < p_1\)}
    {\(p_1 = p_3 < p_2\)}
    {\(p_1 = p_2 = p_3\)}
\end{problem}

\begin{answer}
D
\end{answer}

\begin{solution}
简单随机抽样时，每个个体被抽中的概率为 \(n/N\)；系统抽样的抽样间隔为 \(N/n\)，每个个体被抽中的概率仍为 \(n/N\)；分层抽样按比例抽取时，若第 \(j\) 层有 \(N_j\) 个个体并抽取 \(n_j=nN_j/N\) 个，则该层每个个体被抽中的概率为 \(n_j/N_j=n/N\). 因此 \(p_1=p_2=p_3\)，故选 D.
\end{solution}

\begin{problem}
下列函数中，既是偶函数又在区间 \((- \infty, 0)\) 上单调递增的是（\quad）

\choices
    {\(f(x) = \frac{1}{x^{2}}\)}
    {\(f(x) = x^{2} + 1\)}
    {\(f(x) = x^{3}\)}
    {\(f(x) = 2^{-x}\)}
\end{problem}

\begin{answer}
A
\end{answer}

\begin{solution}
选项 A 中，\(f(-x)=1/(-x)^2=f(x)\)，所以它是偶函数；且在 \((-\infty,0)\) 上 \(f'(x)=-2/x^3>0\)，所以单调递增. 选项 B 在该区间上导数 \(2x<0\)，单调递减；选项 C 是奇函数；选项 D 满足 \(f(-x)=2^x\neq2^{-x}\)（取 \(x\neq0\)），不是偶函数. 因此选 A.
\end{solution}

\begin{problem}
在区间 \([-2,3]\) 上随机选取一个数 \(X\)，则 \(X \leqslant 1\) 的概率为（\quad）

\choices
    {\(\frac{4}{5}\)}
    {\(\frac{3}{5}\)}
    {\( \frac{2}{5} \)}
    {\(\frac{1}{5}\)}
\end{problem}

\begin{answer}
B
\end{answer}

\begin{solution}
在区间 \([-2,3]\) 上随机选取一个数，取值长度为 \(5\). 满足 \(X\leqslant1\) 的取值区间为 \([-2,1]\)，长度为 \(3\)，所以 \(P(X\leqslant1)=3/5\)，故选 B.
\end{solution}

\begin{problem}
若圆 \(C_1: x^2 + y^2 = 1\) 与圆 \(C_2: x^2 + y^2 - 6x - 8y + m = 0\) 外切，则 \(m =\)（\quad）

\choices
    {21}
    {19}
    {9}
    {\(-11\)}
\end{problem}

\begin{answer}
C
\end{answer}

\begin{solution}
将 \(C_2\) 配方得 \((x-3)^2+(y-4)^2=25-m\)，所以两圆心之间的距离为 \(5\)，且 \(C_2\) 的半径为 \(r=\sqrt{25-m}\). 两圆外切时 \(5=1+r\)，故 \(r=4\). 于是 \(25-m=16\)，解得 \(m=9\)，故选 C.
\end{solution}

\begin{problem}
执行如图所示的程序框图，如果输入的 \(t \in [-2, 2]\)，则输出的 \(S\) 属于（\quad）

\texfigure[max width=0.38\linewidth]{img_repaint/2014/hunan_liberal/q7_fig1.tex}

\choices
    {\([-6, - 2]\)}
    {\([-5, - 1]\)}
    {\([-4, 5]\)}
    {\([-3,6]\)}
\end{problem}

\begin{answer}
D
\end{answer}

\begin{solution}
当 \(t\in[0,2]\) 时，程序直接输出 \(S=t-3\)，故 \(S\in[-3,-1]\). 当 \(t\in[-2,0)\) 时，先把 \(t\) 更新为 \(2t^2+1\)，此值为正，随后输出 \(S=2t^2-2\)，故 \(S\in(-2,6]\). 两种情形合并得 \(S\in[-3,6]\)，故选 D.
\end{solution}

\begin{problem}
一块石材表示的几何体的三视图如图所示，将该石材切削，打磨，加工成球，则能得到的最大球的半径等于（\quad）

\bitmapfigure[max width=0.45\linewidth]{img/2014/hunan_liberal/q8_fig1.png}

\choices
    {1}
    {2}
    {3}
    {4}
\end{problem}

\begin{answer}
B
\end{answer}

\begin{solution}
由三视图可知该几何体是底面为直角三角形、长度为 \(12\) 的直三棱柱，直角三角形的两直角边长分别为 \(6\)，\(8\)，斜边长为 \(10\). 球的最大半径不能超过三角形截面的内切圆半径，且不能超过棱柱长度的一半. 三角形内切圆半径为 \(\frac{6+8-10}{2}=2\)，而棱柱长度的一半为 \(6\). 将球心置于棱柱中部并使截面圆为直角三角形的内切圆，可以得到半径为 \(2\) 的球，所以最大半径为 \(2\)，故选 B.
\end{solution}

\begin{problem}
若 \(0 < x_{1} < x_{2} < 1\)，则（\quad）

\choices
    {\(\e^{x_2} - \e^{x_1} > \ln x_2 - \ln x_1\)}
    {\(\e^{x_2} - \e^{x_1} < \ln x_2 - \ln x_1\)}
    {\(x_{2}\e^{x_{1}} > x_{1}\e^{x_{2}}\)}
    {\(x_{2}\e^{x_{1}} < x_{1}\e^{x_{2}}\)}
\end{problem}

\begin{answer}
C
\end{answer}

\begin{solution}
因为 \(0<x_1<x_2<1\)，所以在 \([x_1,x_2]\) 上有 \(1/t>1\)，从而 \(\ln x_2-\ln x_1=\displaystyle\int_{x_1}^{x_2}\frac{1}{t}\,\mathrm dt>x_2-x_1\). 这等价于 \(\ln(x_2/x_1)>x_2-x_1\)，再取指数可得 \(x_2\e^{x_1}>x_1\e^{x_2}\)，故选项 C 正确. 选项 A、B 的比较结果并不固定：当 \(t\in[1/4,1/2]\) 时，\(t\e^t\leqslant \frac12\e^{1/2}<1\)，所以 \(\e^t<1/t\)，取 \(x_1=1/4\)，\(x_2=1/2\) 可得 \(\e^{x_2}-\e^{x_1}<\ln x_2-\ln x_1\)；当 \(t\in[4/5,9/10]\) 时，\(t\e^t\geqslant\frac45(1+\frac45)>1\)，所以 \(\e^t>1/t\)，取 \(x_1=4/5\)，\(x_2=9/10\) 可得相反的不等式. 因此 A、B 均不能恒成立，D 与 C 矛盾.
\end{solution}

\begin{problem}
在平面直角坐标系中，\(O\) 为原点，\(A(-1,0)\)，\(B(0,\sqrt{3})\)，\(C(3,0)\). 动点 \(D\) 满足 \(\left|\overrightarrow{CD}\right| = 1\)，则 \(\left|\overrightarrow{OA} + \overrightarrow{OB} + \overrightarrow{OD}\right|\) 的取值范围是（\quad）

\choices
    {[4, 6]}
    {\([\sqrt{19} - 1, \sqrt{19} + 1]\)}
    {\([2\sqrt{3}, 2\sqrt{7}]\)}
    {\([\sqrt{7} - 1, \sqrt{7} + 1]\)}
\end{problem}

\begin{answer}
D
\end{answer}

\begin{solution}
设 \(\overrightarrow{CD}=\bs u\)，则 \(|\bs u|=1\)，且 \(\overrightarrow{OA}+\overrightarrow{OB}+\overrightarrow{OD}=(-1,0)+(0,\sqrt3)+(3,0)+\bs u=(2,\sqrt3)+\bs u\). 由于 \(|(2,\sqrt3)|=\sqrt7\)，由三角不等式及其逆不等式得 \(\sqrt7-1\leqslant|(2,\sqrt3)+\bs u|\leqslant\sqrt7+1\). 当 \(\bs u\) 与 \((2,\sqrt3)\) 同向或反向时分别取得两个端点，因此取值范围为 \([\sqrt7-1,\sqrt7+1]\)，故选 D.
\end{solution}

\section{填空题}

\begin{problem}
复数 \(\frac{3 + \i}{\i^2}\)（\(\i\) 为虚数单位）的实部等于\fillinblank{}.
\end{problem}

\begin{answer}
\(-3\)
\end{answer}

\begin{solution}
因为 \(\i^2=-1\)，所以 \(\frac{3+\i}{\i^2}=\frac{3+\i}{-1}=-3-\i\). 因此该复数的实部为 \(-3\).
\end{solution}

\begin{problem}
在平面直角坐标系中，曲线 \(C: \begin{cases}
x = 2 + \frac{\sqrt{2}}{2} t,\\
y = 1 + \frac{\sqrt{2}}{2} t,
\end{cases}\) （\(t\) 为参数）的普通方程\fillinblank{}.
\end{problem}

\begin{answer}
\(x-y-1=0\)
\end{answer}

\begin{solution}
由参数方程两式相减，得 \(x-y=1\)，所以普通方程为 \(x-y-1=0\). 反之，令该直线上的点满足 \(x-y=1\)，即可取 \(t=\sqrt2(x-2)=\sqrt2(y-1)\)，故参数方程表示的正是这条直线.
\end{solution}

\begin{problem}
若变量 \(x\)，\(y\) 满足约束条件 \(\begin{cases}
y \leqslant x,\\
x + y \leqslant 4,\\
y \geqslant 1,
\end{cases}\) 则 \(z = 2x + y\) 的最大值为\fillinblank{}.
\end{problem}

\begin{answer}
\(7\)
\end{answer}

\begin{solution}
约束区域由 \(y=1\)，\(y=x\) 和 \(x+y=4\) 围成，其三个顶点分别为 \((1,1)\)，\((3,1)\)，\((2,2)\). 线性函数 \(z=2x+y\) 在多边形上的最大值在顶点处取得，且 \(z(1,1)=3\)，\(z(3,1)=7\)，\(z(2,2)=6\). 因此 \(z\) 的最大值为 \(7\).
\end{solution}

\begin{problem}
平面上一机器人在行进中始终保持与点 \(F(1,0)\) 的距离和到直线 \(x = -1\) 的距离相等. 若机器人接触不到过点 \(P(-1,0)\) 且斜率为 \(k\) 的直线，则 \(k\) 的取值范围是\fillinblank{}.
\end{problem}

\begin{answer}
\((-\infty,-1)\cup(1,+\infty)\)
\end{answer}

\begin{solution}
设机器人位置为 \(Q(x,y)\). 由抛物线的定义，\(QF\) 到焦点 \(F(1,0)\) 的距离等于 \(Q\) 到准线 \(x=-1\) 的距离，因此 \((x-1)^2+y^2=(x+1)^2\)，化简得轨迹方程 \(y^2=4x\). 过 \(P(-1,0)\) 且斜率为 \(k\) 的直线为 \(y=k(x+1)\). 当 \(k\ne0\) 时，代入抛物线得 \(k^2x^2+(2k^2-4)x+k^2=0\)，其判别式为 \(\Delta=16(1-k^2)\). 此时直线与抛物线没有公共点，当且仅当 \(\Delta<0\)，即 \(|k|>1\). 当 \(k=0\) 时，直线为 \(y=0\)，与抛物线交于顶点 \((0,0)\)，不符合题意. 所以 \(k\) 的取值范围为 \((-\infty,-1)\cup(1,+\infty)\).
\end{solution}

\begin{problem}
若 \( f(x) = \ln\left(\e^{3x} + 1\right) + ax \) 是偶函数，则 \( a = \)\fillinblank{}.
\end{problem}

\begin{answer}
\(-\frac{3}{2}\)
\end{answer}

\begin{solution}
由 \(f(x)\) 为偶函数，对任意 \(x\) 有 \(f(x)=f(-x)\). 而 \(f(-x)=\ln(\e^{-3x}+1)-ax=\ln(\e^{3x}+1)-3x-ax\). 与 \(f(x)=\ln(\e^{3x}+1)+ax\) 比较，得 \((2a+3)x=0\) 对任意 \(x\) 成立，因此 \(2a+3=0\)，即 \(a=-\frac32\).
\end{solution}

\section{解答题}

\begin{problem}
已知数列 \(\{a_{n}\}\) 的前 \(n\) 项和 \(S_{n} = \frac{n^{2} + n}{2}\)，\(n \in \N^{*}\).

\begin{enumerate}
\item 求数列 \( \{a_{n}\} \) 的通项公式；
\item 设 \( b_{n} = 2^{a_{n}} + (-1)^{n}a_{n} \)，求数列 \( \{b_{n}\} \) 的前 \( 2n \) 项和.
\end{enumerate}
\end{problem}

\begin{solution}
\begin{enumerate}
\item 当 \(n=1\) 时，\(a_1=S_1=1\). 当 \(n\geqslant2\) 时，\(a_n=S_n-S_{n-1}=\frac{n^2+n}{2}-\frac{(n-1)^2+(n-1)}{2}=n\). 所以 \(a_n=n\)（\(n\in\N^*\)）.
\item 由 \(a_k=k\)，得 \(b_k=2^k+(-1)^k k\). 因此
\[\sum_{k=1}^{2n}b_k=\sum_{k=1}^{2n}2^k+\sum_{k=1}^{2n}(-1)^k k=(2^{2n+1}-2)+[(-1+2)+\cdots+(-(2n-1)+2n)]=2^{2n+1}+n-2\]
\end{enumerate}
\end{solution}

\begin{problem}
某企业有甲，乙两个研发小组，为了比较他们的研发水平，现随机抽取这两个小组往年研发新产品的结果如下：\((a, b)\)，\((a, \overline{b})\)，\((a, b)\)，\((\overline{a}, b)\)，\((\overline{a}, \overline{b})\)，
\((a, b)\)，\((a, b)\)，\((a, \overline{b})\)，\((\overline{a}, b)\)，\((a, \overline{b})\)，
\((\overline{a}, \overline{b})\)，\((a, b)\)，\((a, \overline{b})\)，\((\overline{a}, b)\)，\((a, b)\)，其中 \(a\)，\(\overline{a}\) 分别表示甲组研发成功和失败；\(b\)，\(\overline{b}\) 分别表示乙组研发成功和失败.

\begin{enumerate}
\item 若某组成功研发一种新产品，则给该组记 1 分，否则记 0 分，试计算甲，乙两组研发新产品的成绩的平均数和方差，并比较甲，乙两组的研发水平；
\item 若该企业安排甲，乙两组各自研发一种新产品，试估计恰有一组研发成功的概率.
\end{enumerate}
\end{problem}

\begin{solution}
\begin{enumerate}
\item 统计 15 组结果，甲组成功 10 次、失败 5 次，乙组成功 9 次、失败 6 次. 由于成绩只取 0 和 1，甲组成绩的平均数为 \(\bar X_A=10/15=2/3\)，方差为 \(D(X_A)=\frac{10}{15}\left(1-\frac23\right)^2+\frac{5}{15}\left(0-\frac23\right)^2=\frac29\). 乙组成绩的平均数为 \(\bar X_B=9/15=3/5\)，方差为 \(D(X_B)=\frac{9}{15}\left(1-\frac35\right)^2+\frac{6}{15}\left(0-\frac35\right)^2=\frac{6}{25}\). 因为 \(\bar X_A>\bar X_B\) 且 \(D(X_A)<D(X_B)\)，所以甲组平均成绩较高、成绩更稳定，估计甲组研发水平较高.
\item 恰有一组成功对应于 \((a,\overline b)\) 或 \((\overline a,b)\). 前者出现 4 次，后者出现 3 次，所以用样本频率估计所求概率为 \(\frac{4+3}{15}=\frac{7}{15}\).
\end{enumerate}
\end{solution}

\begin{problem}
如图，已知二面角 \(\alpha - MN - \beta\) 的大小为 \(60^{\circ}\)，菱形 \(ABCD\) 在面 \(\beta\) 内，\(A\)，\(B\) 两点在棱 \(MN\) 上，\(\angle BAD = 60^{\circ}\)，\(E\) 是 \(AB\) 的中点，\(DO \perp\) 面 \(\alpha\)，垂足为 \(O\).

\begin{enumerate}
\item 证明： \(AB \perp\) 平面 \(ODE\)；
\item 求异面直线 \(BC\) 与 \(OD\) 所成角的余弦值.
\end{enumerate}

\bitmapfigure[max width=0.45\linewidth]{img/2014/hunan_liberal/q18_fig1.png}
\end{problem}

\begin{solution}
\begin{enumerate}
\item 因为 \(ABCD\) 是菱形且 \(\angle BAD=60^\circ\)，所以 \(AB=AD\)，从而三角形 \(ABD\) 为等边三角形. \(E\) 是 \(AB\) 的中点，故 \(DE\perp AB\). 又因为 \(AB\subset MN=\alpha\cap\beta\)，且 \(DO\perp\) 平面 \(\alpha\)，所以 \(DO\perp AB\). 由 \(\overrightarrow{OE}=\overrightarrow{OD}+\overrightarrow{DE}\)，得 \(OE\perp AB\). 直线 \(DE\) 与 \(OE\) 在 \(E\) 相交且都在平面 \(ODE\) 内，所以 \(AB\perp\) 平面 \(ODE\).
\item 设 \(AB=s\)，取 \(E\) 为原点，令 \(AB\) 为 \(x\) 轴，平面 \(\alpha\) 为 \(z=0\)，平面 \(\beta\) 与平面 \(\alpha\) 的二面角为 \(60^\circ\). 则可取 \(A=(-s/2,0,0)\)，\(B=(s/2,0,0)\). 由等边三角形 \(ABD\) 得 \(DE=s\sqrt3/2\)，故可取 \(D=(0,s\sqrt3/4,3s/4)\)，从而 \(O=(0,s\sqrt3/4,0)\). 由平行四边形关系，\(\overrightarrow{BC}=\overrightarrow{AD}=(s/2,s\sqrt3/4,3s/4)\)，而 \(\overrightarrow{OD}=(0,0,3s/4)\). 于是
\[\cos\theta=\frac{|\overrightarrow{BC}\cdot\overrightarrow{OD}|}{|BC|\,|OD|}=\frac{9s^2/16}{s\cdot3s/4}=\frac34\]
\end{enumerate}
\end{solution}

\begin{problem}
如图，在平面四边形 \(ABCD\) 中，\(DA \perp AB\)，\(DE = 1\)，\(EC = \sqrt{7}\)，\(EA = 2\)，\(\angle ADC = \frac{2\pi}{3}\)，\(\angle BEC = \frac{\pi}{3}\).

\begin{enumerate}
\item 求 \( \sin \angle CED \) 的值；
\item 求 \(BE\) 的长.
\end{enumerate}

\bitmapfigure[max width=0.38\linewidth]{img/2014/hunan_liberal/q19_fig1.png}
\end{problem}

\begin{solution}
\begin{enumerate}
\item 因为 \(A\)，\(E\)，\(D\) 共线，所以 \(\angle EDC=\angle ADC=2\pi/3\). 在三角形 \(DEC\) 中，设 \(DC=d\)，由余弦定理得 \(7=1+d^2-2\cdot1\cdot d\cos(2\pi/3)=1+d^2+d\). 解得 \(d=2\). 于是 \(\cos\angle CED=\frac{EC^2+ED^2-CD^2}{2EC\cdot ED}=\frac{7+1-4}{2\sqrt7}=\frac{2}{\sqrt7}\)，所以 \(\sin\angle CED=\sqrt{1-4/7}=\frac{\sqrt{21}}{7}\).
\item 设 \(\theta=\angle CED\). 由于 \(EA\) 与 \(ED\) 反向，\(\angle AEC=\pi-\theta\). 在图示凸四边形中，射线 \(EB\) 位于 \(\angle AEC\) 内，故 \(\angle AEB=\angle AEC-\angle BEC=2\pi/3-\theta\). 由上问得 \(\cos\theta=2/\sqrt7\)，\(\sin\theta=\sqrt{21}/7\)，从而 \(\cos\angle AEB=\cos(2\pi/3-\theta)=-\frac12\cdot\frac{2}{\sqrt7}+\frac{\sqrt3}{2}\cdot\frac{\sqrt{21}}7=\frac{\sqrt7}{14}\). 又因为 \(AE\perp AB\)，三角形 \(AEB\) 在 \(A\) 处为直角，故 \(\cos\angle AEB=AE/BE=2/BE\)，于是 \(BE=28/\sqrt7=4\sqrt7\).
\end{enumerate}
\end{solution}

\begin{problem}
如图，\(O\) 为坐标原点，双曲线 \(C_{1}:\frac{x^{2}}{a_{1}^{2}} - \frac{y^{2}}{b_{1}^{2}} = 1\)（\(a_{1} > 0\)，\(b_{1} > 0\)） 和椭圆 \(C_{2}:\frac{y^{2}}{a_{2}^{2}} + \frac{x^{2}}{b_{2}^{2}} = 1\)（\(a_{2} > b_{2} > 0\)） 均过点 \(P\left(\frac{2\sqrt{3}}{3},1\right)\)，且以 \(C_{1}\) 的两个顶点和 \(C_{2}\) 的两个焦点为顶点的四边形是面积为 \(2\) 的正方形.

\begin{enumerate}
\item 求 \(C_1\)，\(C_2\) 的方程；
\item 是否存在直线 \(l\)，使得 \(l\) 与 \(C_1\) 交于 \(A\)，\(B\) 两点，与 \(C_2\) 只有一个公共点，且 \(\left|\overrightarrow{OA} + \overrightarrow{OB}\right| = \left|\overrightarrow{AB}\right|\)，证明你的结论.
\end{enumerate}

\bitmapfigure[max width=0.42\linewidth]{img/2014/hunan_liberal/q20_fig1.png}
\end{problem}

\begin{solution}
\begin{enumerate}
\item 令 \(c_2=\sqrt{a_2^2-b_2^2}\)，则 \(C_1\) 的两个顶点为 \((\pm a_1,0)\)，\(C_2\) 的两个焦点为 \((0,\pm c_2)\). 四点构成的正方形的两条对角线长分别为 \(2a_1\)，\(2c_2\)，所以 \(a_1=c_2\)，且正方形面积为 \(2a_1c_2=2\). 由此得 \(a_1=c_2=1\). 点 \(P\) 在 \(C_1\) 上，故 \(\frac{4}{3a_1^2}-\frac1{b_1^2}=1\)，解得 \(b_1^2=3\). 点 \(P\) 在 \(C_2\) 上，且 \(a_2^2-b_2^2=c_2^2=1\)，故 \(\frac1{a_2^2}+\frac{4}{3b_2^2}=1\). 令 \(u=b_2^2\)，则 \(a_2^2=u+1\)，代入得 \(\frac1{u+1}+\frac4{3u}=1\)，即 \(3u^2-4u-4=0\). 因 \(u>0\)，得 \(u=2\)，从而 \(a_2^2=3\). 所以 \(C_1:\frac{x^2}{1}-\frac{y^2}{3}=1\)，\(C_2:\frac{y^2}{3}+\frac{x^2}{2}=1\).
\item 先考虑非竖直直线 \(l:y=kx+m\). 代入 \(C_1\) 得 \((3-k^2)x^2-2kmx-(m^2+3)=0\). 因为 \(l\) 与 \(C_1\) 交于两个不同点，所以 \(3-k^2\ne0\)；当 \(3-k^2=0\) 时上式至多为一次方程，不可能有两个交点. 设两个交点的横坐标为 \(x_A\)，\(x_B\)，由韦达定理及根的差得到 \(\overrightarrow{OA}+\overrightarrow{OB}=\left(\frac{2km}{3-k^2},\frac{6m}{3-k^2}\right)\)，以及 \(|AB|^2=\frac{12(1+k^2)(m^2+3-k^2)}{(3-k^2)^2}\). 条件 \(|\overrightarrow{OA}+\overrightarrow{OB}|=|AB|\) 化平方并约去非零因子 \(3-k^2\)，得 \(m^2=\frac32(1+k^2)\). 另一方面，直线与椭圆 \(x^2/2+y^2/3=1\) 只有一个公共点，等价于代入后关于 \(x\) 的二次方程判别式为零，即 \(m^2=3+2k^2\). 两式合并得到 \(\frac32(1+k^2)=3+2k^2\)，即 \(k^2=-3\)，不可能.

竖直直线设为 \(x=h\). 它与 \(C_1\) 的两个交点为 \((h,\pm\sqrt{3(h^2-1)})\)，条件 \(|\overrightarrow{OA}+\overrightarrow{OB}|=|AB|\) 给出 \(2|h|=2\sqrt{3(h^2-1)}\)，即 \(h^2=3/2\). 但它与椭圆相切时必须为椭圆的左右端切线，故 \(h^2=2\)，矛盾. 因此不存在满足条件的直线 \(l\).
\end{enumerate}
\end{solution}

\begin{problem}
已知函数 \( f(x) = x \cos x - \sin x + 1\)（\(x > 0\)）.

\begin{enumerate}
\item 求 \( f(x) \) 的单调区间；
\item 记 \(x_{i}\) 为 \(f(x)\) 的从小到大的第 \(i\)（\(i \in \N^{*}\)）个零点，证明：对一切 \(n \in \N^{*}\)，有 \(\frac{1}{x_1^2} + \frac{1}{x_2^2} + \cdots + \frac{1}{x_n^2} < \frac{2}{3}\).
\end{enumerate}
\end{problem}

\begin{solution}
\begin{enumerate}
\item \(f'(x)=\cos x-x\sin x-\cos x=-x\sin x\). 因为 \(x>0\)，所以在 \(((2k)\pi,(2k+1)\pi)\) 上 \(\sin x>0\)，函数单调递减；在 \(((2k+1)\pi,(2k+2)\pi)\) 上 \(\sin x<0\)，函数单调递增，其中 \(k\in\N\).
\item 连续延拓定义 \(f(0)=1\). 在每个区间 \(((2k)\pi,(2k+1)\pi)\) 上，有 \(f(2k\pi)=2k\pi+1>0\)，\(f((2k+1)\pi)=1-(2k+1)\pi<0\)，函数从正值减到负值，恰有一个零点；在每个区间 \(((2k+1)\pi,(2k+2)\pi)\) 上，这两个端点的函数值分别为负值和正值，函数单调递增，也恰有一个零点. 特别地，\(f((4k+1)\pi/2)=0\)，所以这些递减区间内的零点为 \(x_{2k+1}=(4k+1)\pi/2\)；其余零点满足 \(x_{2k+2}>(2k+1)\pi\). 因此对所有 \(i\geqslant2\)，有 \(x_i>(i-1)\pi\)，且 \(x_1=\pi/2\).

当 \(n=1\) 时，\(1/x_1^2=4/\pi^2<2/3\). 当 \(n\geqslant2\) 时，利用 \(x_i>(i-1)\pi\) 得 \(\displaystyle\sum_{i=1}^{n}\frac1{x_i^2}<\frac4{\pi^2}+\frac1{\pi^2}\sum_{j=1}^{n-1}\frac1{j^2}\). 又因为对 \(j\geqslant2\) 有 \(1/j^2<1/[j(j-1)]\)，所以 \(\sum_{j=1}^{n-1}1/j^2<2\). 于是 \(\displaystyle\sum_{i=1}^{n}\frac1{x_i^2}<\frac6{\pi^2}<\frac6{3^2}=\frac23\)（因为 \(\pi>3\)），证毕.
\end{enumerate}
\end{solution}
